\section{引言：从语言模型到序列建模}

本讲是 CS224N 第六讲，承接上一讲（Lecture 5）关于语言模型与循环神经网络的基础内容，进一步深入讨论三个核心主题：（1）语言模型的评估方法——困惑度（Perplexity）；（2）RNN 的梯度问题与 LSTM 的解决方案；（3）序列到序列模型在机器翻译中的应用。

\begin{figure}[H]
\centering
\includegraphics[width=0.95\textwidth]{slides-images/slide-000.jpg}
\caption{CS224N Lecture 6 封面：LSTM RNNs and Neural Machine Translation\protect\footnotemark}
\end{figure}
\footnotetext{Slides 第1页。}

\begin{knowledgebox}{前置知识回顾}
上一讲（Lecture 5）介绍了两个独立但相关的概念：
\begin{itemize}
    \item \textbf{语言模型}（Language Model）：预测下一个词的系统，是 NLP 中许多任务（文本生成、似然估计）的基础组件
    \item \textbf{循环神经网络}（RNN）：能处理任意长度序列输入的神经网络架构，在每个时间步应用相同的权重，可选择性地在每步产生输出
\end{itemize}
这两个概念常常结合使用，但 RNN 也可用于语言模型以外的其他序列任务。
\end{knowledgebox}

\subsection{课程结构概览}

本讲的内容组织如下：

\begin{enumerate}
    \item \textbf{困惑度}（Perplexity）：语言模型的标准评估指标
    \item \textbf{RNN 的梯度问题}：梯度消失与梯度爆炸
    \item \textbf{LSTM}：通过门控机制解决长距离依赖问题
    \item \textbf{RNN 的其他应用}：双向 RNN、多层 RNN
    \item \textbf{机器翻译}：从统计方法到神经方法的革命
    \item \textbf{Seq2Seq 模型}：编码器--解码器架构
\end{enumerate}

\subsection{本章小结}

本讲从语言模型的评估出发，经过 RNN 的梯度问题分析和 LSTM 的设计动机，最终引入序列到序列模型在机器翻译中的应用。这条主线贯穿了从理论分析到工程实践的完整链路。

%% ========== 困惑度 ==========

